% SGA 6, Exposición XIV: traducción española alineada con la fuente.
% Alcance exacto: fr.tex líneas 15726--15825, secciones 5 y 6 completas.
% Las lecturas fuente corregidas se conservan mediante \sisi.

\providecommand{\sisi}[2]{#2}

\section*{5. Relaciones entre $\operatorname{Gr}^\bullet(X)$ y
$H^{2\bullet}(X,\Zell(\bullet))$}

\textbf{5.1.} El ejemplo dado en 4.8 no excluye que quizá pueda encontrarse,
para un esquema arbitrario $X$, un homomorfismo de anillos
\[
g^\bullet:\operatorname{Gr}^\bullet(X)\longrightarrow
H^{2\bullet}(X,\Zell(\bullet)),
\]
tal que, para todo módulo localmente libre $E$ sobre $X$, $g^i$ transforme
la clase de Chern $c_i(E)$ de la Exp. V en la clase de Chern cohomológica de
SGA 5 VII. Naturalmente, $\ell$ denota aquí un número primo distinto de las
características residuales de $X$. Como el anillo
$\operatorname{Gr}^\bullet(X)$ está engendrado por los $\sisi{c^i}{c_i}(E)$, se ve en % Canon SGA-SRC-001421
todo caso que $g^\bullet$ está determinado de manera única por la condición de ser
compatible con las clases de Chern.

\textbf{5.2.} Utilizando el homomorfismo de Chern
\[
K^\bullet(X)\longrightarrow
\Ch\bigl(H^{2\bullet}(X,\Zell(\bullet))\bigr),
\]
se obtiene, pasando a los graduados asociados respecto de la filtración del
primer miembro y de la filtración por el grado del segundo, un homomorfismo
de grupos graduados, pero no de anillos,
\[
\psi:\operatorname{Gr}^\bullet(X)\longrightarrow
H^{2\bullet}(X,\Zell(\bullet)),
\]
que satisface la relación
\[
\psi(c_i(E))=(-1)^{i-1}(i-1)!\,c_i^\ell(E),
\]
donde $c_i^\ell(E)$ denota la clase de Chern $\ell$-ádica. Se deduce de los
$\psi^i$, mediante división, un homomorfismo de anillos
\[
g_\Q:\operatorname{Gr}^\bullet(X)\longrightarrow
H^{2\bullet}(X,\Qell(\bullet))=
H^{2\bullet}(X,\Zell(\bullet))\otimes_{\Zell}\Qell,
\]
que transforma clases de Chern en clases de Chern y cuya imagen del primer
miembro está contenida en la de
$\sisi{H^2(X,\Zell(\bullet))}{H^{2\bullet}(X,\Zell(\bullet))}$. Por tanto,
la cuestión de la existencia de $g^\bullet$ se resuelve afirmativamente si se
reemplaza $H^{2\bullet}(X,\Zell(\bullet))$ por su cociente por el subgrupo de
torsión. El primer caso dudoso al definir las componentes $g^i$ del
hipotético homomorfismo $g$ es el de $g^3$; a priori solo se consigue definir
\[
2g^3=\psi^3.
\]
La respuesta a la cuestión de la existencia de $g^3$ sería negativa si, por
ejemplo, se pudieran encontrar tres haces invertibles $L_1,L_2,L_3$ sobre
$X$ que satisfagan la relación
\[
(L_1-1)(L_2-1)(L_3-1)=0\quad\text{en }K^\bullet(X),
\]
pero tales que
\[
c_1^\ell(L_1)c_1^\ell(L_2)c_1^\ell(L_3)\ne0
\quad\text{en }H^6(X,\mu_2^{\otimes3}).
\]

\textbf{5.3.} Cuando $X$ es liso y cuasiproyectivo sobre un cuerpo $k$,
puede plantearse análogamente la existencia de un homomorfismo de anillos
\[
g:\operatorname{Gr}^\bullet(X)\longrightarrow A(X)
\]
con valores en el anillo de Chow $A(X)$ y que transforme clases de Chern en
clases de Chern. Se pueden repetir al respecto las consideraciones
precedentes, que muestran en particular cómo construir directamente
$(-1)^{i-1}(i-1)!g^i$ cuando no se dispone del propio $g^i$.

\textbf{5.4.} Señalemos, para terminar este número, que se plantea
esencialmente la misma cuestión si, en lugar de un esquema $X$, se considera
un espacio topológico compacto, digamos, $X$. Todavía se puede considerar el
graduado $\operatorname{Gr}^\bullet(X)$ de $K^\bullet(X)$, provisto de su
$\lambda$-filtración, y buscar un homomorfismo de anillos
\[
\operatorname{Gr}^\bullet(X)\longrightarrow H^{2\bullet}(X,\Z)
\]
que transforme clases de Chern en clases de Chern.

\section*{6. Teorema de Riemann--Roch cohomológico y homomorfismo de Gysin}

\textbf{6.1.} Para simplificar, nos limitaremos aquí a esquemas $X$ que
admiten un módulo invertible amplio, lo que evita distinguir entre el
$K'(X)$ ingenuo y el otro. Comencemos por la fórmula de Riemann--Roch «con
denominadores» para un morfismo propio de intersección completa $f:X\to Y$.
Como, para un número primo $\ell$ distinto de las características residuales
de $Y$, se pueden definir las clases de Chern $\ell$-ádicas de un módulo
localmente libre $E$,
\[
c_i(E)\in H^{2i}(X,\Zell(i))
\]
(SGA 5 VII), correspondientes a un homomorfismo
\begin{equation}\tag{6.1}
K^\bullet(X)\longrightarrow\Ch\bigl(H^{2\bullet}(X,\Zell(\bullet))\bigr),
\end{equation}
se plantea la cuestión de dar una variante cohomológica de la fórmula de
Riemann--Roch. El único ingrediente que falta para dar sentido a la fórmula
es un homomorfismo de Gysin
\begin{equation}\tag{6.2}
H^n(X,\Qell(i))\longrightarrow H^{n-2d}(Y,\Qell(i-d)),
\end{equation}
donde $d$ es la dimensión relativa virtual de $f$. Suponiendo definido este
homomorfismo y teniendo en cuenta la fórmula de Riemann--Roch en la forma de
la Exp. VIII, con valores en $\operatorname{Gr}^\bullet(Y)_\Q$, se comprueba que la
validez de la fórmula de Riemann--Roch $\ell$-ádica para $f$ equivale
exactamente a la conmutatividad del cuadrado
\begin{equation}\tag{6.3}
\begin{tikzcd}[column sep=large,row sep=large]
\operatorname{Gr}^\bullet(X)_{\Q} \arrow[r] \arrow[d] &
H^{2\bullet}(X,\Qell(\bullet)) \arrow[d] \\
\operatorname{Gr}^\bullet(Y)_{\Q} \arrow[r] &
H^{2\bullet}(Y,\Qell(\bullet)),
\end{tikzcd}
\end{equation}
donde las flechas horizontales se definen a partir de (6.1), pasando a los
graduados asociados y dividiendo de modo que transformen clases de Chern en
clases de Chern, y las flechas verticales son, respectivamente, la flecha de
la Exp. VIII y el homomorfismo de Gysin en cohomología, supuesto dado.

\textbf{6.2.} Obsérvese que el homomorfismo de Gysin está definido en todo
caso si $Y$ es liso sobre un cuerpo $k$ (SGA 5 IV), utilizando el teorema de
pureza cohomológica relativa sobre $k$. La misma construcción se aplica
cuando $Y$ es excelente, regular y de característica cero, gracias al
teorema de pureza cohomológica absoluta de SGA 4 XIX, y se aplicará
automáticamente a todo $Y$ excelente y regular, sin restricción de
característica, una vez demostrado en toda generalidad el teorema de pureza
cohomológica absoluta. Finalmente, la misma construcción se aplica, gracias
al teorema de pureza cohomológica relativa (SGA 4 XVI 3.7), cuando $X$ e $Y$
son ambos lisos sobre un mismo esquema $S$ y $f$ es un $S$-morfismo. Cabe
conjeturar que, en cada uno de estos casos en que está definido el
homomorfismo de Gysin, el diagrama (6.3) es conmutativo. El único caso en que
esta afirmación está actualmente verificada es aquel en que $X$ e $Y$ son
ambos lisos sobre un cuerpo $k$; resulta entonces de la conmutatividad del
diagrama análogo a (6.3)
\begin{equation}\tag{6.4}
\begin{tikzcd}[column sep=large,row sep=large]
A(X) \arrow[r] \arrow[d] & H^{2\bullet}(X,\Zell(\bullet)) \arrow[d] \\
A(Y) \arrow[r] & H^{2\bullet}(Y,\Zell(\bullet)),
\end{tikzcd}
\end{equation}
en el que intervienen los anillos de clases de ciclos módulo equivalencia
racional $A(X)$ y $A(Y)$, cuya conmutatividad se establece en SGA 5 IV.

\textbf{6.3.} Obsérvese que, en cada uno de los casos considerados arriba en
los que se consigue definir un homomorfismo de Gysin, se obtiene en realidad
una construcción para la cohomología con coeficientes en $\Zell(\bullet)$,
y no solo en $\Qell(\bullet)$. Esto permite, cuando $f$ es una inmersión,
plantear la validez de la fórmula de Riemann--Roch «sin denominadores», como
se explica en el número 3. Cuando $X$ e $Y$ son ambos lisos sobre un cuerpo
$k$, la conmutatividad de (6.4) muestra que bastaría demostrar la fórmula sin
denominadores en el anillo de Chow $A(Y)$; cf. (3.3).

\textbf{6.4.} Cuando no se impone ninguna hipótesis de regularidad sobre $Y$,
e incluso si $X$ e $Y$ son de tipo finito sobre el cuerpo $\C$, ya no se ve
una definición plausible de un homomorfismo de Gysin cohomológico (6.2), ni
siquiera para una inmersión regular de codimensión $1$. ¿Puede encontrarse,
sin embargo, uno con propiedades razonables, en particular que haga
conmutativo el cuadrado (6.3)? La cuestión está evidentemente relacionada
con la de definir, para un $X$ de intersección completa de codimensión $i$
en $Y$, «su clase de cohomología»
\[
\gamma(X)\in H^{2i}(Y,\Zell(i)),
\]
de manera que transforme intersecciones sin componentes excedentarias en
cup-productos. Esta última cuestión se resuelve afirmativamente, al menos
módulo torsión, es decir, en $H^{2\bullet}(Y,\Qell(\bullet))$, definiendo
\[
\gamma(X)=\frac{1}{(-1)^{i-1}(i-1)!}\,c_i(\mathcal O_X).
\]
Pero ignoro si $c_i(\mathcal O_X)$ es divisible por $(i-1)!$ en
$H^{2i}(Y,\Zell(i))$. En particular, si $Y$ es, por ejemplo, un esquema
proyectivo de dimensión $3$ sobre el cuerpo de los complejos y $X$ se reduce
a un punto singular de $Y$, ¿es $c_3(\mathcal O_X)$ divisible por $2$ en
$H^6(Y,\Z_2(3))$ y, en particular, es nula la clase de Chern en
$H^6(Y,\Z/2\Z)$? Señalemos también, en la misma dirección, la cuestión
siguiente: sea $Y$ un esquema proyectivo singular de dimensión $4$, sea $X$
un subesquema cerrado de dimensión $1$ que sea globalmente una intersección
de tres divisores positivos $D_1,D_2,D_3$, y consideremos el producto de sus
clases de Chern módulo $2$:
\[
c_1(D_1)c_1(D_2)c_1(D_3)\in H^6(Y,\Z/2\Z).
\]
¿Es este elemento de $H^6$ independiente de la elección particular de los
$D_i$ cuya intersección es $X$?
